% figures15.tex -- Chapter 15 (Analytic Geometry) figures redrawn in TikZ.
% Coordinate-geometry chapter: every figure is exactly constructible from the
% mathematics in the text (axes, plotted points, lines, circles), so vertices
% are computed, never eyeballed. Proportions were measured off the printed
% figures at 400 dpi (PDF pp. 267-279) during the figure review; the readings
% are recorded in chapters/ch15-UNCERTAIN.md.
%
% Helper macros used only inside this file (check_labels greps FIG[A-Z]+ for
% the figure macros, so these helpers are not picked up as figures).
% `positioning' is needed for the "above left=<a> and <b>" node syntax; it is
% loaded here rather than in brumfiel.sty because this file is also compiled
% standalone by tools/check_labels.py.
\usetikzlibrary{positioning}

% Standard axes WITH the letters x and y: {xmin}{xmax}{ymin}{ymax}.
% The book letters the axes in Figs 15-6 to 15-19 and leaves them unlettered
% in Fig 15-4, so the two variants are kept apart.
\newcommand{\XVaxes}[4]{%
  \draw[fig,-latex] (#1,0) -- (#2,0);
  \node[vlab,right] at (#2,0) {$x$};
  \draw[fig,-latex] (0,#3) -- (0,#4);
  \node[vlab,right] at (0,#4) {$y$};}
% Same axes, unlettered (Fig 15-4).
\newcommand{\XVaxesnl}[4]{%
  \draw[fig,-latex] (#1,0) -- (#2,0);
  \draw[fig,-latex] (0,#3) -- (0,#4);}
% Origin label
\newcommand{\XVzero}{\node[vlab,below left=-1pt and 0pt] at (0,0) {$0$};}
% Unit ticks on the x-axis, from #1 to #2 step 1 (skipping 0)
\newcommand{\XVxticks}[2]{%
  \foreach \t in {#1,...,#2}{\ifnum\t=0\else
    \draw[tick] (\t,-0.11) -- (\t,0.11);\fi}}
\newcommand{\XVyticks}[2]{%
  \foreach \t in {#1,...,#2}{\ifnum\t=0\else
    \draw[tick] (-0.11,\t) -- (0.11,\t);\fi}}

% ---------------------------------------------------------------- Fig 15-1
% Coordinatized line: O, the unit point P_1, a positive point P_x and a
% negative point P'_{x'}.  Collinear by construction.  Spacings follow the
% printed figure (O -> P_1 : P_1 -> P_x : P'_{x'} -> O = 1 : 1.55 : 1.22).
\newcommand{\FIGXVONE}{\begin{bkfigure}[\linewidth]
\begin{tikzpicture}[scale=1.0]
  \draw[fig] (-0.3,0) -- (10.3,0);
  \foreach \x/\below/\above in {0.6/{P'_{x'}}/{x'}, 3.5/{O}/{},
                                5.9/{P_1}/{1}, 9.6/{P_x}/{x}}
    {\dt{\x,0}
     \node[vlab,below=2pt] at (\x,0) {$\below$};}
  \foreach \x/\above in {0.6/{x'}, 5.9/{1}, 9.6/{x}}
    {\node[vlab,above=2pt] at (\x,0) {$\above$};}
\end{tikzpicture}
\figcap{15--1}
\end{bkfigure}}

% ---------------------------------------------------------------- Fig 15-2
% Two PERPENDICULAR lines through O, coordinatized, drawn obliquely to make
% the point that the choice of axes is arbitrary.
% abscissa positive direction 210 deg; ordinate positive direction 300 deg.
% Measured off the plate: abscissa runs 2.6 units back / 5.0 forward, ordinate
% 1.5 back / 3.4 forward.  Every numeral sits on the side the book puts it on:
% the positive abscissa numerals on the 120 deg side, the positive ordinate
% numerals on the 30 deg side, and each -1 thrown to the far side of its own
% axis so the two do not crowd the 0.
\newcommand{\FIGXVTWO}{\begin{bkfigure}[\linewidth]
\begin{tikzpicture}[scale=1.040]
  % unit vectors
  \def\ax{-0.86603} \def\ay{-0.5}      % abscissa positive (210 deg)
  \def\ox{0.5}      \def\oy{-0.86603}  % ordinate positive (300 deg)
  % the two axes
  \draw[fig] (2.6*-\ax,2.6*-\ay) -- (5.0*\ax,5.0*\ay);
  \draw[fig] (1.5*-\ox,1.5*-\oy) -- (3.4*\ox,3.4*\oy);
  % ticks + numerals on the abscissa axis (tick direction = ordinate dir)
  \foreach \t/\n in {1/1, 2/2, 3/3}
    {\draw[tick] ({\t*\ax-0.20*\ox},{\t*\ay-0.20*\oy})
              -- ({\t*\ax+0.20*\ox},{\t*\ay+0.20*\oy});
     \node[slab] at ({\t*\ax-0.46*\ox},{\t*\ay-0.46*\oy}) {$\n$};}
  \draw[tick] ({-\ax-0.20*\ox},{-\ay-0.20*\oy})
           -- ({-\ax+0.20*\ox},{-\ay+0.20*\oy});
  \node[slab] at ({-\ax-0.50*\ox},{-\ay-0.50*\oy}) {$-1$};
  % ticks + numerals on the ordinate axis (tick direction = abscissa dir)
  \foreach \t/\n in {1/1, 2/2, 3/3}
    {\draw[tick] ({\t*\ox-0.20*\ax},{\t*\oy-0.20*\ay})
              -- ({\t*\ox+0.20*\ax},{\t*\oy+0.20*\ay});
     \node[slab] at ({\t*\ox-0.50*\ax},{\t*\oy-0.50*\ay}) {$\n$};}
  \draw[tick] ({-\ox-0.20*\ax},{-\oy-0.20*\ay})
           -- ({-\ox+0.20*\ax},{-\oy+0.20*\ay});
  \node[slab] at ({-\ox+0.66*\ax},{-\oy+0.66*\ay}) {$-1$};
  % 0 in the open sector between the two negative rays, as in the book
  \node[slab] at (0.16,0.60) {$0$};
  \node[slab,left]  at (5.15*\ax,5.15*\ay) {Abscissas};
  \node[slab,right] at (3.55*\ox,3.55*\oy) {Ordinates};
\end{tikzpicture}
\figcap{15--2}
\end{bkfigure}}

% ---------------------------------------------------------------- Fig 15-3
% Theorem 15-1: P,Q in the SAME quadrant => PQ meets neither axis;
% P,Q in DIFFERENT quadrants => PQ meets at least one axis.
\newcommand{\FIGXVTHREE}{\begin{bkfigure}[\linewidth]
\begin{tikzpicture}[scale=1.014]
  \draw[fig] (-4.3,0) -- (3.9,0);
  \draw[fig] (0,2.9) -- (0,-1.5);
  % same quadrant (I): segment meets neither axis
  \coordinate (P1) at (1.6,2.15);  \coordinate (Q1) at (2.65,1.35);
  \draw[key] (P1) -- (Q1);
  \node[vlab,above=1pt] at (P1) {$P$};
  \node[vlab,right=1pt] at (Q1) {$Q$};
  % quadrant II to quadrant III: crosses the axis of abscissas
  \coordinate (Q2) at (-2.45,1.5); \coordinate (P2) at (-3.75,-0.62);
  \draw[key] (Q2) -- (P2);
  \node[vlab,above=1pt] at (Q2) {$Q$};
  \node[vlab,left=1pt]  at (P2) {$P$};
  % quadrant II to quadrant IV: crosses both axes
  \coordinate (P3) at (-1.25,1.05); \coordinate (Q3) at (2.35,-0.72);
  \draw[key] (P3) -- (Q3);
  \node[vlab,above=1pt] at (P3) {$P$};
  \node[vlab,right=1pt] at (Q3) {$Q$};
  \foreach \p in {P1,Q1,Q2,P2,P3,Q3} {\dt{\p}}
\end{tikzpicture}
\figcap{15--3}
\end{bkfigure}}

% ---------------------------------------------------------------- Fig 15-4
% The four quadrants, numbered.  The book leaves this pair of axes unlettered
% (no x, no y) and runs the left branch longer than the right.
\newcommand{\FIGXVFOUR}{\begin{bkfigure}[\linewidth]
\begin{tikzpicture}[scale=0.936]
  \XVaxesnl{-3.6}{2.9}{-1.2}{2.7}
  \XVxticks{1}{2}
  \XVyticks{1}{2}
  \node[slab,above] at (1,0.12) {$1$};
  \node[slab,above] at (2,0.12) {$2$};
  \node[slab,left]  at (-0.12,1) {$1$};
  \node[slab,left]  at (-0.12,2) {$2$};
  \XVzero
  \node[vlab] at (1.40,1.50)  {\textup{I}};
  \node[vlab] at (-1.70,1.50) {\textup{II}};
  \node[vlab] at (-1.75,-0.87){\textup{III}};
  \node[vlab] at (1.36,-0.87) {\textup{IV}};
\end{tikzpicture}
\figcap{15--4}
\end{bkfigure}}

% ---------------------------------------------------------------- Fig 15-5
% Rotated (still perpendicular) axes; the reader is asked which quadrant is
% which.  Ordinate positive direction 150 deg, abscissa positive 240 deg --
% the printed pair measures about 92 deg apart, so they are squared up.
% Branch lengths follow the plate: the ordinate axis runs 1.8 units forward
% and 3.9 back, the abscissa 2.05 forward and 1.95 back.
\newcommand{\FIGXVFIVE}{\begin{bkfigure}[\linewidth]
\begin{tikzpicture}[scale=1.040]
  \def\ox{-0.86603} \def\oy{0.5}       % ordinate positive (150 deg)
  \def\ax{-0.5}     \def\ay{-0.86603}  % abscissa positive (240 deg)
  \draw[fig,-latex] (3.9*-\ox,3.9*-\oy) -- (1.8*\ox,1.8*\oy);
  \draw[fig,-latex] (1.95*-\ax,1.95*-\ay) -- (2.05*\ax,2.05*\ay);
  \dt{\ox,\oy}
  \dt{\ax,\ay}
  \node[slab,above] at (\ox,\oy) {$1$};
  \node[slab,right] at (\ax,\ay) {$1$};
  \node[slab,above] at (2.0*\ox,2.0*\oy) {Ordinates};
  \node[slab,below] at (2.25*\ax,2.25*\ay) {Abscissas};
\end{tikzpicture}
\figcap{15--5}
\end{bkfigure}}

% ---------------------------------------------------------------- Fig 15-6
% Definition 15-4: l_1 through P parallel to the axis of ordinates meets the
% first axis in A; l_2 through P parallel to the axis of abscissas meets the
% second axis in B.  A sits on the positive abscissa axis and B on the
% negative ordinate axis, so P falls in quadrant IV, as in the book.
\newcommand{\FIGXVSIX}{\begin{bkfigure}[\linewidth]
\begin{tikzpicture}[scale=1.040]
  \XVaxes{-1.6}{2.8}{-3.2}{1.65}
  \coordinate (A) at (1.55,0);
  \coordinate (B) at (0,-1.70);
  \coordinate (P) at (1.55,-1.70);
  \draw[aux] (1.55,1.55) -- (1.55,-3.5);
  \draw[aux] (-1.65,-1.70) -- (3.2,-1.70);
  \foreach \p in {A,B,P} {\dt{\p}}
  \node[vlab,above left=-1pt and 0pt] at (0,0) {$O$};
  \node[vlab,above right=1pt and 1pt] at (A) {$A$};
  \node[vlab,above left=1pt and 1pt]  at (B) {$B$};
  \node[vlab,above right=1pt and 1pt] at (P) {$P$};
  \node[vlab,right] at (3.2,-1.70) {$l_2$};
  \node[vlab,below right=-1pt and 0pt] at (1.55,-3.5) {$l_1$};
\end{tikzpicture}
\figcap{15--6}
\end{bkfigure}}

% ---------------------------------------------------------------- Fig 15-7
% The line x = 3: parallel to the y-axis, 3 units to its right.  The book
% ticks only 1 and 2 -- the line itself stands in for the third mark -- and
% sets the label to the RIGHT of the line.
\newcommand{\FIGXVSEVEN}{\begin{bkfigure}[\linewidth]
\begin{tikzpicture}[scale=0.910]
  \XVaxes{-0.9}{4.3}{-1.7}{3.55}
  \XVxticks{1}{2}
  \draw[key] (3,-1.65) -- (3,3.48);
  \XVzero
  \node[vlab,right] at (3.08,1.2) {$x=3$};
\end{tikzpicture}
\figcap{15--7}
\end{bkfigure}}

% ------------------------------------------------------------ Fig 15-8, 15-9
% 15-8: the points of 2x - 3y + 6 = 0 found in Example 1, namely
%       (0,2), (3,4), (-3,0) -- collinear by construction.  Both axes are
%       ticked in the book.
% 15-9: the line y = -2 of Example 2, parallel to the x-axis, 2 units below.
%       The axis of ordinates stops on that line, as it does in the book.
\newcommand{\FIGXVEIGHTNINE}{\begin{bkfigure}[\linewidth]
\begin{minipage}{0.48\linewidth}\centering
\begin{tikzpicture}[scale=0.44]
  \XVaxes{-4.6}{6.4}{-2.6}{5.4}
  \XVxticks{-4}{5}
  \XVyticks{-2}{4}
  % the line 2x - 3y + 6 = 0, i.e. y = (2/3)x + 2 -- endpoints computed
  \draw[aux] (-4.2,-0.8) -- (5.4,5.6);
  \foreach \p in {(-3,0),(0,2),(3,4)} {\node[bkdot] at \p {};}
\end{tikzpicture}
\figcap{15--8}
\end{minipage}\hfill
\begin{minipage}{0.48\linewidth}\centering
\begin{tikzpicture}[scale=0.44]
  \XVaxes{-3.4}{5.3}{-2.05}{3.1}
  \draw[aux] (-3.6,-2) -- (5.4,-2);
  \draw[tick] (-0.22,-1) -- (0.22,-1);
  \XVzero
  \node[vlab,above=1pt] at (3.15,-2) {$(x,-2)$};
\end{tikzpicture}
\figcap{15--9}
\end{minipage}
\end{bkfigure}}

% ---------------------------------------------------------- Fig 15-10, 15-11
% 15-10: proof of Theorem 15-2, Case 3.  (0,-C/B), (x1,y1), (x2,y2) are
%        collinear; the two vertical legs are perpendicular to the base.
%        The base is taken at height 0 and the x-axis sits below it, at the
%        same fraction of x2 as in the book.
% 15-11: proof of Theorem 15-3, Case 3.  (x1,y1), (x2,y2), (x,y) collinear;
%        right angles at (x2,y1) and (x,y1).  In the book (x1,y1) lies just
%        to the LEFT of the axis of ordinates and the whole construction sits
%        above the axis of abscissas.
\newcommand{\FIGXVTENELEVEN}{\begin{bkfigure}[\linewidth]
\begin{minipage}{0.49\linewidth}\centering
\begin{tikzpicture}[scale=0.55]
  % base line y = -C/B taken at height 0; the x-axis sits below it.
  \def\bY{0}
  \coordinate (O)  at (0,\bY);          % (0, -C/B)
  \coordinate (P1) at (2.81,1.47525);   % (x1, y1)   slope 0.525
  \coordinate (P2) at (4.2,2.205);      % (x2, y2)   -- collinear with O,P1
  \coordinate (F1) at (2.81,\bY);
  \coordinate (F2) at (4.2,\bY);
  \draw[fig,-latex] (-0.52,-1.65) -- (4.54,-1.65);
  \node[vlab,right] at (4.54,-1.65) {$x$};
  \draw[fig,-latex] (0,-2.20) -- (0,2.49);
  \node[vlab,right] at (0,2.49) {$y$};
  \draw[fig] (O) -- (F2);
  \draw[fig] (P1) -- (F1);
  \draw[fig] (P2) -- (F2);
  \draw[aux] (O) -- (P2);
  \foreach \p in {O,P1,P2} {\dt{\p}}
  \node[slab,left]  at (O)  {$(0,-C/B)$};
  \node[slab,anchor=south east,xshift=-1pt,yshift=3pt] at (P1) {$(x_1,y_1)$};
  \node[slab,anchor=south east,xshift=9pt,yshift=3pt]  at (P2) {$(x_2,y_2)$};
  % one node holds both foot labels, as the book sets them: side by side on a
  % single line below the base.  Two separate nodes overlap at this trim.
  \node[slab,below] at (3.3,\bY) {$(x_1,-C/B)\quad(x_2,-C/B)$};
\end{tikzpicture}
\figcap{15--10}
\end{minipage}\hfill
\begin{minipage}{0.49\linewidth}\centering
\begin{tikzpicture}[scale=0.60]
  \coordinate (P1) at (-0.61,1.36);     % (x1, y1)
  \coordinate (P2) at (1.78,1.77825);   % (x2, y2)   slope 0.175
  \coordinate (P3) at (4.67,2.284);     % (x, y)   -- collinear with P1,P2
  \coordinate (F2) at (1.78,1.36);
  \coordinate (F3) at (4.67,1.36);
  \draw[fig,-latex] (-1.65,0) -- (4.93,0);
  \node[vlab,right] at (4.93,0) {$x$};
  \draw[fig,-latex] (0,-1.02) -- (0,3.44);
  \node[vlab,right] at (0,3.44) {$y$};
  \draw[fig] (P1) -- (F3);
  \draw[fig] (P2) -- (F2);
  \draw[fig] (P3) -- (F3);
  \draw[fig] (P1) -- (P3);
  \foreach \p in {P1,P2,P3} {\dt{\p}}
  \node[slab,left] at (P1) {$(x_1,y_1)$};
  \node[slab] at ($(P2)+(-0.34,0.48)$) {$(x_2,y_2)$};
  \node[slab] at ($(P3)+(0.10,0.50)$)  {$(x,y)$};
  \node[slab,below] at (F2) {$(x_2,y_1)$};
  \node[slab,below] at (F3) {$(x,y_1)$};
\end{tikzpicture}
\figcap{15--11}
\end{minipage}
\end{bkfigure}}

% --------------------------------------------------------------- Fig 15-12
% The line through (-2,3) and (1,5) of the worked Example.  Both axes ticked;
% the book sets (-2,3) above-left of its point and (1,5) below-right of its
% own, which also keeps both clear of the line.
\newcommand{\FIGXVTWELVE}{\begin{bkfigure}[\linewidth]
\begin{tikzpicture}[scale=0.676]
  \XVaxes{-5.0}{5.6}{-1.2}{6.9}
  \XVxticks{-4}{5}
  \XVyticks{1}{5}
  % slope 2/3 through (-2,3) and (1,5) -- endpoints computed, not eyeballed
  \draw[key] (-4.1,1.6) -- (3.45,6.63333);
  \dt{-2,3}
  \dt{1,5}
  \XVzero
  \node[vlab,above left=0pt and 0pt]  at (-2,3) {$(-2,3)$};
  \node[vlab,below right=0pt and 0pt] at (1,5)  {$(1,5)$};
\end{tikzpicture}
\figcap{15--12}
\end{bkfigure}}

% ---------------------------------------------------------- Fig 15-13, 15-14
% 15-13: distances along lines parallel to an axis.  AB = 1 - (-3) = 4;
%        CD = |x2 - x1|.  All four dashed pairs are axis-parallel.  C and D
%        straddle the axis of ordinates, at (-1,-1) and (2,-1), as in the
%        book, and each of the four lettered points carries its letter on one
%        side of the dot and its coordinate pair on the other.
% 15-14: the right triangle P1 P2 Q of the distance formula, right angle at Q.
\newcommand{\FIGXVTHIRTEENFOURTEEN}{\begin{bkfigure}[\linewidth]
\begin{minipage}{0.58\linewidth}\centering
\begin{tikzpicture}[scale=0.52]
  \draw[fig,-latex] (-2.8,0) -- (7.0,0);
  \node[vlab,right] at (7.0,0) {$x$};
  \draw[fig,-latex] (0,-3.6) -- (0,4.6);
  \node[vlab,right] at (0,4.6) {$y$};
  \node[slab,above left=-1pt and 0pt] at (0,0) {$0$};
  % horizontal pair (2,3)-(6,3)
  \draw[aux] (2,3) -- (6,3);
  \dt{2,3} \dt{6,3}
  \node[slab,above] at (2,3.1) {$(2,3)$};
  \node[slab,above] at (6,3.1) {$(6,3)$};
  % horizontal pair (-2,2)-(3,2)
  \draw[aux] (-2,2) -- (3,2);
  \dt{-2,2} \dt{3,2}
  \node[slab,left]  at (-2.1,2) {$(-2,2)$};
  \node[slab,right] at (3.1,2)  {$(3,2)$};
  % vertical pair A(1,1)-B(1,-3)
  \draw[aux] (1,1) -- (1,-3);
  \dt{1,1} \dt{1,-3}
  \node[slab,left]  at (0.88,1)  {$A$};
  \node[slab,right] at (1.12,1)  {$(1,1)$};
  \node[slab,left]  at (0.88,-3) {$B$};
  \node[slab,right] at (1.12,-3) {$(1,-3)$};
  % horizontal pair C(x1,y1)-D(x2,y1), straddling the axis of ordinates.
  % Dropped to y = -1.5: at this trim a label set above a point at y = -1
  % would run back over the axis of abscissas.
  \draw[aux] (-1,-1.5) -- (2,-1.5);
  \dt{-1,-1.5} \dt{2,-1.5}
  \node[slab,above] at (-1,-1.4) {$C$};
  \node[slab,above] at (2,-1.4)  {$D$};
  \node[slab,left]  at (-1.12,-1.5) {$(x_1,y_1)$};
  \node[slab,right] at (2.12,-1.5)  {$(x_2,y_1)$};
\end{tikzpicture}
\figcap{15--13}
\end{minipage}\hfill
\begin{minipage}{0.38\linewidth}\centering
\begin{tikzpicture}[scale=0.62]
  \draw[fig,-latex] (-0.95,0) -- (4.57,0);
  \node[vlab,right] at (4.57,0) {$x$};
  \draw[fig,-latex] (0,-0.26) -- (0,3.55);
  \node[vlab,right] at (0,3.55) {$y$};
  \node[slab,below right=-1pt and 1pt] at (0,0) {$0$};
  \coordinate (P1) at (0.62,1.15);
  \coordinate (P2) at (3.72,3.05);
  \coordinate (Q)  at (3.72,1.15);
  \draw[key] (P1) -- (P2);
  \draw[aux] (P1) -- (Q);
  \draw[aux] (Q)  -- (P2);
  \foreach \p in {P1,P2,Q} {\dt{\p}}
  \node[vlab,above,xshift=-2pt] at (P1) {$P_1$};
  \node[vlab,anchor=south east,xshift=-1pt] at (P2) {$P_2$};
  \node[vlab,right,yshift=5pt] at (Q) {$Q$};
  % nudged right so it clears the axis of ordinates, as the book's does
  \node[slab,below,xshift=8pt] at (P1) {$(x_1,y_1)$};
  % the book prints (x_2,y_1) at P_2 as well as at Q -- see ch15-UNCERTAIN.md
  \node[slab,anchor=south west,xshift=1pt] at (P2) {$(x_2,y_1)$};
  \node[slab,below=3pt] at (Q) {$(x_2,y_1)$};
\end{tikzpicture}
\figcap{15--14}
\end{minipage}
\end{bkfigure}}

% ---------------------------------------------------------- Fig 15-15, 15-16
% 15-15: M is the midpoint of P1P2; parallels to the axes through M bisect
%        the sides of the right triangle.
% 15-16: the worked example -- M(0,1) is the midpoint of (-1,3) and (1,-1).
\newcommand{\FIGXVFIFTEENSIXTEEN}{\begin{bkfigure}[\linewidth]
\begin{minipage}{0.52\linewidth}\centering
\begin{tikzpicture}[scale=0.756]
  \draw[fig,-latex] (-0.40,0) -- (3.60,0);
  \node[vlab,right] at (3.60,0) {$x$};
  \draw[fig,-latex] (0,-0.35) -- (0,3.90);
  \node[vlab,right] at (0,3.90) {$y$};
  \node[slab,below right=-1pt and 1pt] at (0,0) {$0$};
  \coordinate (P1) at (0.46,1.05);
  \coordinate (P2) at (3.36,2.52);
  \coordinate (M)  at ($(P1)!0.5!(P2)$);
  \coordinate (C)  at (3.36,1.05);
  \coordinate (Mh) at (3.36,1.785);    % foot of the parallel to the x-axis
  \coordinate (Mv) at (1.91,1.05);     % foot of the parallel to the y-axis
  \draw[key] (P1) -- (P2);
  \draw[aux] (P1) -- (C) -- (P2);
  \draw[aux] (M) -- (Mh);
  \draw[aux] (M) -- (Mv);
  \foreach \p in {P1,P2,M} {\dt{\p}}
  \node[slab,anchor=north west,xshift=-9pt,yshift=-3pt] at (P1) {$P_1\,(x_1,y_1)$};
  \node[slab,right] at (P2) {$P_2\,(x_2,y_2)$};
  \node[vlab,above left=-3pt and -1pt] at (M) {$M$};
  % Solidus form: a built-up fraction draws a rule inside the label box, which
  % the collision audit reads as geometry.  The solidus runs about two and a
  % half times as wide as the book's stacked fraction, so the label is lifted
  % clear of P_2 rather than tucked beside it.  See chapters/ch15-UNCERTAIN.md.
  \node[slab,anchor=south west] at (0.05,2.95)
    {$\left((x_1+x_2)/2,\ (y_1+y_2)/2\right)$};
\end{tikzpicture}
\figcap{15--15}
\end{minipage}\hfill
\begin{minipage}{0.44\linewidth}\centering
\begin{tikzpicture}[scale=0.651]
  \XVaxes{-2.0}{3.2}{-1.95}{3.5}
  \node[slab,below left=-1pt and 0pt] at (0,0) {$0$};
  \dt{-1,3}
  \dt{0,1}
  \dt{1,-1}
  \node[slab,left]  at (-1.15,3) {$P_1\,(-1,3)$};
  \node[slab,right] at (0.18,1)  {$M\,(0,1)$};
  \node[slab,right] at (1.18,-1) {$P_2\,(1,-1)$};
\end{tikzpicture}
\figcap{15--16}
\end{minipage}
\end{bkfigure}}

% --------------------------------------------------------------- Fig 15-17
% Exercise *7: the segment joining the midpoints of two sides is parallel to
% the third side and half as long.  Vertices (0,0), (a,0), (b,c); the book
% draws b barely larger than a, so the right side is all but vertical, and
% the midline is a plain dashed segment with no arrowhead.
\newcommand{\FIGXVSEVENTEEN}{\begin{bkfigure}[\linewidth]
\begin{tikzpicture}[scale=0.936]
  \draw[fig,-latex] (-0.70,0) -- (3.70,0);
  \node[vlab,right] at (3.70,0) {$x$};
  \draw[fig,-latex] (0,-0.60) -- (0,2.75);
  \node[vlab,right] at (0,2.75) {$y$};
  \coordinate (A) at (0,0);
  \coordinate (B) at (3.10,0);
  \coordinate (C) at (3.15,2.48);
  \draw[fig] (A) -- (B) -- (C) -- cycle;
  \coordinate (M1) at ($(A)!0.5!(C)$);
  \coordinate (M2) at ($(B)!0.5!(C)$);
  \draw[aux] (M1) -- (M2);
  \foreach \p in {M1,M2} {\dt{\p}}
  \node[slab,below right=-1pt and 0pt] at (A) {$(0,0)$};
  \node[slab,below] at (B) {$(a,0)$};
  \node[slab,right] at (C) {$(b,c)$};
\end{tikzpicture}
\figcap{15--17}
\end{bkfigure}}

% ---------------------------------------------------------- Fig 15-18, 15-19
% 15-18: the circle of Eq. (15-6): centre (h,k), radius r, a point (x,y) on it.
%        The book's centre sits just above the axis of abscissas and about
%        half a radius right of the axis of ordinates, so the circle crosses
%        both axes; the axis of ordinates stops well short of the bottom of
%        the circle.
% 15-19: the circle x^2+y^2+6x-8y+19 = 0, centre (-3,4), radius sqrt(6).
\newcommand{\FIGXVEIGHTEENNINETEEN}{\begin{bkfigure}[\linewidth]
\begin{minipage}{0.50\linewidth}\centering
\begin{tikzpicture}[scale=0.66]
  \draw[fig,-latex] (-1.82,0) -- (4.90,0);
  \node[vlab,right] at (4.90,0) {$x$};
  \draw[fig,-latex] (0,-0.77) -- (0,2.65);
  \node[vlab,right] at (0,2.65) {$y$};
  % the book sets 0 below left of the origin; the origin here lies inside the
  % circle, and the audit approximates each quarter-arc by its chord, so a
  % label dropped below left reads as touching that chord.  Set above left.
  \node[slab,above left=-1pt and 0pt] at (0,0) {$0$};
  \coordinate (K) at (1.04,0.38);
  \def\rr{1.85}
  \draw[key] (K) circle (\rr);
  % (x,y) placed on the circle at 148.5 deg from the centre -- exact by
  % construction, so the drawn radius really has length r.
  \coordinate (X) at ($(K)+(148.5:\rr)$);
  \draw[fig] (K) -- (X);
  \dt{K}
  \node[slab,right] at (K) {$(h,k)$};
  \node[slab,left=1pt]  at (X) {$(x,y)$};
  % r set off the radius along the perpendicular, so it does not sit on the line
  \node[slab] at ($(K)!0.5!(X)+(58.5:0.34)$) {$r$};
\end{tikzpicture}
\figcap{15--18}
\end{minipage}\hfill
\begin{minipage}{0.46\linewidth}\centering
\begin{tikzpicture}[scale=0.48]
  \draw[fig,-latex] (-6.0,0) -- (3.2,0);
  \node[vlab,right] at (3.2,0) {$x$};
  \draw[fig,-latex] (0,-1.0) -- (0,7.0);
  \node[vlab,right] at (0,7.0) {$y$};
  \node[slab,below right=-1pt and 1pt] at (0,0) {$0$};
  \XVxticks{-5}{-1}
  \XVyticks{1}{5}
  \draw[key] (-3,4) circle (2.4494897);
  \dt{-3,4}
  \node[slab,below=3pt] at (-3,4) {$(-3,4)$};
\end{tikzpicture}
\figcap{15--19}
\end{minipage}
\end{bkfigure}}
